\documentclass[11pt]{article}
\usepackage[margin=1in]{geometry}
\usepackage{enumitem}
% =============================================================================
% Preambles/header.tex — shared LaTeX/Beamer preamble for this project
%
% Use from a Beamer lecture by prepending:
%     \input{header}
% and compiling with TEXINPUTS=../Preambles:$TEXINPUTS (the /compile-latex
% skill does this automatically).
%
% PALETTE CONTRACT. Color names defined here MUST match the names in
% Quarto/theme-template.scss so Beamer and Quarto renderings use the same
% palette. The scripts/check-palette-sync.sh script enforces this.
%
% To customize: edit the HEX values below AND the matching $variable in
% Quarto/theme-template.scss, then run scripts/check-palette-sync.sh.
% =============================================================================

% -----------------------------------------------------------------------------
% Engine detection + encoding
%
% Our /compile-latex skill uses XeLaTeX, where inputenc is unnecessary and
% can produce encoding warnings. Only load inputenc under pdfTeX.
% -----------------------------------------------------------------------------
\usepackage{iftex}
\ifPDFTeX
  \usepackage[utf8]{inputenc}
\fi

\usepackage{amsmath, amssymb, amsthm}
\usepackage{booktabs}
\usepackage{graphicx}

% -----------------------------------------------------------------------------
% PALETTE — mirrors Quarto/theme-template.scss variable names.
% Load xcolor and define colors BEFORE anything that references them
% (notably hyperref, hypersetup, and the Beamer color assignments).
% -----------------------------------------------------------------------------
\usepackage{xcolor}

\definecolor{primary-blue}{HTML}{012169}      % headings, accents
\definecolor{primary-gold}{HTML}{B9975B}      % emphasis, borders
\definecolor{highlight-yellow}{HTML}{F2A900}  % markers, alerts
\definecolor{light-bg}{HTML}{E8EDF5}          % title / section backgrounds
\definecolor{jet}{HTML}{1A1A1A}               % body text

% Semantic colors (used in diagrams and annotations)
\definecolor{positive}{HTML}{15803D}          % good / observed / identified
\definecolor{negative}{HTML}{B91C1C}          % bad / problematic / bias
\definecolor{neutral}{HTML}{525252}           % reference / context

% Highlight accents (match .hi-* classes in the SCSS)
\definecolor{hi-slate}{HTML}{314F4F}
\definecolor{hi-green}{HTML}{2E8B57}
\definecolor{hi-red}{HTML}{C41E3A}

% Semantic aliases so skill templates and snippets can write
% \textcolor{accent}{...} without hard-coding "primary-blue".
\colorlet{accent}{primary-blue}
\colorlet{accent2}{primary-gold}

% -----------------------------------------------------------------------------
% Hyperref — loaded AFTER xcolor + palette so link colors resolve.
% -----------------------------------------------------------------------------
\usepackage{hyperref}
\hypersetup{colorlinks=true, linkcolor=primary-blue, urlcolor=primary-blue,
            citecolor=primary-blue, pdfborder={0 0 0}}

% -----------------------------------------------------------------------------
% Course code — set once per deck (after \input{header}, before \begin{document})
% via \coursecode{CS401}. Shown in the footer on every slide so a deck's course
% is identifiable without opening the title page. Empty by default.
% -----------------------------------------------------------------------------
\makeatletter
\newcommand{\coursecode}[1]{\gdef\@coursecodetext{#1}}
\gdef\@coursecodetext{}
\makeatother

% -----------------------------------------------------------------------------
% Beamer theme assignments (only apply under Beamer).
%
% We detect Beamer via \@ifclassloaded{beamer}, which is the documented,
% non-internal way. This must be wrapped in \makeatletter/\makeatother
% because \@ifclassloaded contains an @-catcode character.
% -----------------------------------------------------------------------------
\makeatletter
\@ifclassloaded{beamer}{%
  \usefonttheme{professionalfonts}
  \setbeamercolor{structure}{fg=primary-blue}
  \setbeamercolor{frametitle}{fg=primary-blue}
  \setbeamercolor{title}{fg=primary-blue}
  \setbeamercolor{subtitle}{fg=primary-gold}
  \setbeamercolor{author}{fg=primary-blue}
  \setbeamercolor{institute}{fg=neutral}
  \setbeamercolor{date}{fg=primary-gold}
  \setbeamercolor{itemize item}{fg=primary-blue}
  \setbeamercolor{itemize subitem}{fg=primary-gold}
  \setbeamercolor{alerted text}{fg=primary-gold}
  \setbeamercolor{block title}{fg=white, bg=primary-blue}
  \setbeamercolor{block body}{bg=light-bg}
  % Semantic block variants: block=definition/concept (blue), exampleblock=
  % worked example (green/positive), alertblock=key takeaway or caution (gold).
  % Keep to <=2 colored boxes per slide across ALL three variants (INV-7).
  \setbeamercolor{block title example}{fg=white, bg=positive}
  \setbeamercolor{block body example}{bg=light-bg}
  \setbeamercolor{block title alerted}{fg=white, bg=primary-gold}
  \setbeamercolor{block body alerted}{bg=light-bg}

  % Minimal footer: course code (left) + page X / N (right), no navigation clutter
  \setbeamertemplate{navigation symbols}{}
  \setbeamertemplate{footline}{%
    \color{neutral}\footnotesize\hspace{0.7em}\@coursecodetext\hfill
    \insertframenumber/\inserttotalframenumber\hspace{0.7em}\vspace{0.3em}}
}{}
\makeatother

% -----------------------------------------------------------------------------
% TikZ setup — libraries the snippet gallery uses
% -----------------------------------------------------------------------------
\usepackage{tikz}
\usetikzlibrary{arrows.meta, positioning, calc, decorations.pathreplacing,
                fit, shapes.geometric, backgrounds}

% Shared TikZ styles — snippets and hand-written diagrams can pick these up
% via `\begin{tikzpicture}[every node/.style={...}]` or by naming specific
% styles (e.g., `\node[dag-node] (X) at (0,0) {X};`).
\tikzset{
  % Node styles for causal DAGs and similar
  dag-node/.style = {
    draw=primary-blue, fill=light-bg, rounded corners=2pt,
    minimum width=1.2cm, minimum height=0.9cm, align=center, font=\small
  },
  decision-node/.style = {
    draw=primary-gold, fill=white, diamond, aspect=1.8,
    minimum width=1.8cm, minimum height=1.2cm, align=center, font=\small,
    inner sep=1pt
  },
  % Edge styles — observed vs counterfactual
  observed-edge/.style = {-{Latex[length=2.5mm]}, thick, draw=jet},
  counterfactual-edge/.style = {-{Latex[length=2.5mm]}, thick, draw=neutral, dashed},
  confound-edge/.style = {-{Latex[length=2.5mm]}, thick, draw=negative, dashed},
  % Data-point styles for plots
  observed-dot/.style = {fill=primary-blue, draw=primary-blue, circle, inner sep=1.2pt},
  counterfactual-dot/.style = {fill=white, draw=neutral, circle, inner sep=1.2pt}
}

% -----------------------------------------------------------------------------
% Convenience macros
% -----------------------------------------------------------------------------
\newcommand{\muted}[1]{\textcolor{neutral}{#1}}
\newcommand{\key}[1]{\textcolor{primary-gold}{\textbf{#1}}}
\newcommand{\good}[1]{\textcolor{positive}{#1}}
\newcommand{\bad}[1]{\textcolor{negative}{#1}}

% Standout frame macro: full-bleed dark background for section transitions.
% Beamer-only (uses \begin{frame}/\setbeamercolor/\usebeamerfont) -- do not
% call from an article-class file (e.g. Notes/<CODE>/*-notes.tex). \muted,
% \key, \good, \bad, and \coursecode above are plain xcolor/gdef and are
% safe to use from either class; verified when Notes/article-class support
% was added.
% Expand: \transitionslide{Your transition title}
\newcommand{\transitionslide}[1]{%
  \begingroup
  \setbeamercolor{background canvas}{bg=primary-blue}
  \begin{frame}[plain]
    \centering
    \vfill
    {\usebeamerfont{title}\color{white}\Large #1\par}
    \vfill
  \end{frame}
  \endgroup
}

% Section divider frame (Beamer-only), an alternative to \transitionslide with a
% darker two-line style: near-black (jet) background, a small white label line
% above a larger gold title, and a thin gold rule along the bottom edge.
% Expand: \sectiondivider{Section 2}{Pointers}
\newcommand{\sectiondivider}[2]{%
  \begingroup
  \setbeamercolor{background canvas}{bg=jet}
  \begin{frame}[plain]
    \vspace*{3.0cm}
    {\color{white}\ttfamily\large #1\par}
    \vspace{0.5cm}
    {\color{primary-gold}\LARGE #2\par}
    \begin{tikzpicture}[remember picture, overlay]
      \draw[primary-gold, line width=2.5pt]
        ($(current page.south west)+(0,0.1)$) -- ($(current page.south east)+(0,0.1)$);
    \end{tikzpicture}
  \end{frame}
  \endgroup
}

% -----------------------------------------------------------------------------
% Channel watermark — "Concepts That Click" (YouTube).
%
% Article-class files (CompetitiveExam Q/A sets, Notes, Assignments) get a
% light diagonal draft watermark on every page plus a centered footer naming
% the channel. Beamer decks get a subtle TikZ background watermark instead
% (the footline already carries the course code). Centralized here so every
% MCQ PDF is branded without per-file edits.
% -----------------------------------------------------------------------------
\makeatletter
\@ifclassloaded{article}{%
  \usepackage{draftwatermark}
  \SetWatermarkText{Concepts That Click}
  \SetWatermarkScale{1}
  \SetWatermarkFontSize{2cm}
  \SetWatermarkLightness{0.86}
  \SetWatermarkAngle{45}
  \usepackage{fancyhdr}
  \pagestyle{fancy}
  \fancyhf{}
  \fancyfoot[C]{\footnotesize\textcolor{neutral}{Concepts That Click~|~YouTube: \href{https://www.youtube.com/@concepts.thatclick}{@concepts.thatclick}}}
  \renewcommand{\headrulewidth}{0pt}
  \renewcommand{\footrulewidth}{0pt}
  \fancypagestyle{plain}{%
    \fancyhf{}%
    \fancyfoot[C]{\footnotesize\textcolor{neutral}{Concepts That Click~|~YouTube: \href{https://www.youtube.com/@concepts.thatclick}{@concepts.thatclick}}}%
    \renewcommand{\headrulewidth}{0pt}%
    \renewcommand{\footrulewidth}{0pt}%
  }
}{}
\@ifclassloaded{beamer}{%
  \setbeamertemplate{background}{%
    \begin{tikzpicture}[remember picture, overlay]
      \node[rotate=45, opacity=0.055, align=center]
        at (current page.center)
        {\Huge\textcolor{neutral}{Concepts That Click}};
    \end{tikzpicture}%
  }
}{}
\makeatother 
\coursecode{JKSSB}
\renewcommand{\thesection}{23.\arabic{section}}
\newtheorem{example}{Example}[section]
\newenvironment{definitionbox}[1]{%
  \par\noindent\textbf{Definition (#1).}\ \itshape
}{\par\vspace{0.3em}}
\title{MS Word High-Yield --- Detailed Lecture Notes}
\author{JKSSB Computer Awareness}
\date{\today}

\begin{document}
\maketitle

\section{Why these Word tools are high-yield}
MS Word is a large program, but JKSSB does not test it evenly. A small set of
review and editing tools is asked again and again, usually as an exact shortcut
or as a near-neighbour distinction. This lesson covers that set: Find and
Replace, spell check, hyperlinks and bookmarks, styles, comments, Track Changes,
Page Break versus Section Break, mail merge, and saving, printing and exporting
to PDF. The aim is not to learn every menu; it is to know each tool's shortcut
and the one thing it is \emph{not}.

\begin{definitionbox}{The high-yield set}
The Word features that competitive-exam items concentrate on: Find/Replace,
spell check and grammar tools, hyperlinks and bookmarks, styles, comments,
Track Changes, Section Break, mail merge, and PDF export.
\end{definitionbox}

\begin{center}
\begin{tabular}{p{0.28\textwidth}p{0.62\textwidth}}
\toprule
\textbf{Term} & \textbf{Meaning in this lesson} \\
\midrule
Find/Replace & The tool to locate text and swap it for other text. \\
Spell check & Word's proofing tool for spelling and grammar. \\
Bookmark & A named location or selection inside a document. \\
Style & A named set of formatting applied to text or a paragraph. \\
Track Changes & A mode that records every edit with its author. \\
Section Break & A break that starts a new section with its own layout. \\
Mail merge & Combining a standard letter with a list of records. \\
PDF & Portable Document Format --- a fixed-layout file. \\
\bottomrule
\end{tabular}
\end{center}

\section{Finding and replacing text}
Word has two related tools. \textbf{Find} locates text in the document and is
opened with \textbf{Ctrl+F}. \textbf{Find and Replace} both locates text and
substitutes new text for it, and is opened with \textbf{Ctrl+H}. Inside the
Find and Replace dialog, \emph{Replace} changes the current match while
\emph{Replace All} changes every match in one step. The same dialog can find
formatting, special characters (such as a paragraph mark) and text that uses a
particular style.

The common confusion is between the two shortcuts. Ctrl+F only finds; it does
not change anything. Ctrl+H is the one that replaces. An item that asks how to
``find and replace'' text expects Ctrl+H.

\begin{example}
A student has typed the word ``recieve'' in ten places and wants ``receive''
everywhere. Press Ctrl+H, type \texttt{recieve} in the Find box and
\texttt{receive} in the Replace box, then choose Replace All. Using Ctrl+F
alone would only highlight the mistakes one by one, without correcting them.
\end{example}

\section{Spelling and grammar tools}
Press \textbf{F7} to start the \textbf{spelling and grammar} check. Word
flags a possible misspelling with a red squiggle and a possible grammar problem
with another underline. Modern Word gathers these checks in the
\textbf{Editor} pane, which groups Spelling, Grammar and Refinements. From the
check you can \emph{Ignore} a flagged word once, \emph{Ignore All} every
occurrence, or \emph{Add} the word to the dictionary so it is accepted from
then on.

Two automatic tools support proofing. \textbf{AutoCorrect} fixes common typing
errors as you type (for example, capitalising the first letter of a sentence),
and the \textbf{thesaurus} (Shift+F7) offers synonyms. The exam usually tests
only the main shortcut, F7. Do not confuse F7 with the PowerPoint slide-show
keys: F5 starts a slide show and F9 has other roles in PowerPoint, but in Word
F7 is proofing.

\begin{definitionbox}{Spell check}
The Word tool, opened with F7, that checks spelling and grammar and offers
corrections, ignoring or adding words to the dictionary.
\end{definitionbox}

\section{Hyperlinks and bookmarks}
A \textbf{hyperlink} is text or an object that jumps to another place when
clicked. Press \textbf{Ctrl+K} to insert one. A hyperlink can point to a web
address, to a file, or to a location inside the same document. When it points
inside the document, it usually points to a \textbf{bookmark}.

A \textbf{bookmark} gives a location or a selection a name, so that text,
a hyperlink or a cross-reference can jump to it. Bookmarks are not printed and
are not visible to the reader by default; they are an internal target. The
clean distinction for the exam is that a hyperlink is the \emph{clickable link},
while a bookmark is the \emph{named target} it points to.

\section{Styles}
A \textbf{style} is a named collection of formatting settings: font, size,
colour, spacing, alignment and indent. Instead of formatting each paragraph by
hand, you apply a style. Applying a \textbf{Heading} style does more than
change the look: it gives the paragraph an outline level. That is why the
\textbf{Navigation Pane} can list the document's headings and why a
\textbf{Table of Contents} can be built and refreshed automatically.

Style-based formatting is consistent and easy to update. If you edit a style,
every paragraph that uses that style changes at once. This is the main
advantage of styles over direct formatting.

\begin{example}
A report uses the Heading 1 style for its chapter titles. Later the author
wants all chapter titles in a larger font. Instead of selecting every title,
the author edits the Heading 1 style once, and every chapter title updates
automatically. The Navigation Pane and the table of contents also stay correct.
\end{example}

\section{Comments and Track Changes}
A \textbf{comment} is a note attached to a point in the document. It lets a
reviewer ask a question or suggest a change \emph{without editing the text
itself}. Comments can be replied to and marked resolved, and they are stored as
metadata shown in the margin or the Reviewing Pane.

\textbf{Track Changes} is different: it records the actual edits. When it is
on, each insertion, deletion and formatting change is marked, together with the
\textbf{author's name} and the time. This is called revision metadata. The
reviewer then \textbf{Accepts} or \textbf{Rejects} each change individually or
all at once. Accepting a change keeps it permanently and, with it, removes the
revision metadata for that change. Rejecting a change removes the edit and
keeps the original text.

\begin{definitionbox}{Track Changes}
A Word mode that records every insertion, deletion and formatting edit along
with the author and time, so a reviewer can accept or reject each change.
\end{definitionbox}

\section{Page Break versus Section Break}
A \textbf{Page Break} pushes the following text onto a new page. It does not
start a new section, so the page keeps the same layout, margins, headers,
footers and page numbering as the rest of the document. A quick way to insert a
page break is \textbf{Ctrl+Enter}.

A \textbf{Section Break} divides the document into independent
\textbf{sections}. Because each section has its own layout settings, a Section
Break is what allows \emph{different headers and footers}, different page
numbering, different margins or a different page orientation within one
document. Section Breaks are inserted from the \textbf{Layout} tab, under
\textbf{Breaks}.

This is one of the most tested Word traps. A Page Break changes only the page; a
Section Break changes the section. Only a Section Break can give a document two
different headers or footers. (TRAP-13)

\begin{example}
An item states, ``Using Page Break allows you to have different headers and
footers in the same document.'' This is wrong. A Page Break keeps the same
section formatting. It is a \textbf{Section Break} that allows different
headers and footers in one document. (PYQ-13)
\end{example}

\section{Mail merge}
\textbf{Mail merge} is the process that combines a \textbf{main document} (a
standard letter or template) with a \textbf{data source} (a list of records
such as names and addresses). The main document contains \textbf{merge fields}
as placeholders, and each record fills those fields to produce one personalised
copy. A common data source is a table of names and addresses held in Excel or
in Outlook contacts.

The typical steps are: create or select the main document; select recipients
(the data source); insert merge fields; preview the results; and finish the
merge by printing the letters or sending them by email. The exam definition to
remember is that mail merge \emph{combines letters with a database of
addresses}. (PYQ-37, PYQ-47)

\begin{definitionbox}{Mail merge}
The Word process that combines a standard document with a data source of
records, producing one personalised copy of the document for each record.
\end{definitionbox}

\section{Saving, printing and exporting to PDF}
Press \textbf{Ctrl+S} to save the document; Word normally saves it in the
\textbf{DOCX} format. Press \textbf{Ctrl+P} to open the print options.
\textbf{Save As} can change the file name, the location and the file format.
Using \textbf{Export} or \textbf{Save As PDF} creates a \textbf{PDF} file.

\textbf{PDF} stands for \textbf{Portable Document Format}. A PDF is a
\emph{fixed-layout} file: it looks the same on every device and is good for
sharing and printing, but it is not meant for easy editing. This pairs with the
separate lesson on file formats, where the extension (for example,
\texttt{.pdf}) and the format (PDF) are kept distinct.

\section{Exam retrieval and common confusions}
Five distinctions prevent most errors on this topic.
\begin{enumerate}
  \item \textbf{Ctrl+F} finds; \textbf{Ctrl+H} finds and replaces. Do not swap
    them.
  \item \textbf{F7} is spelling and grammar in Word; F5 starts a PowerPoint
    slide show. (TRAP-32)
  \item A \textbf{Section Break}, not a Page Break, enables different
    headers/footers in one document. (TRAP-13)
  \item A \textbf{bookmark} is a named target; a \textbf{hyperlink}
    (\textbf{Ctrl+K}) is the clickable link.
  \item \textbf{Mail Merge} combines a letter with an address database; it is
    not Track Changes or AutoCorrect. (TRAP-29)
\end{enumerate}

\begin{example}
An item asks which feature combines letters with a database of addresses.
Trace the meaning: a standard letter plus a list of records is mail merge. If
the item instead asks what records every edit with its author, the answer is
Track Changes, and if it asks what fixes typing errors as you type, the answer
is AutoCorrect. (PYQ-37, PYQ-47)
\end{example}

\subsection*{Exercises}
\begin{enumerate}[label=\arabic*.]
  \item State the shortcuts for Find, Find and Replace, spelling and grammar,
    and inserting a hyperlink.
  \item Explain the difference between a Page Break and a Section Break, and
    say which one allows different headers and footers.
  \item Distinguish a bookmark from a hyperlink.
  \item What does Track Changes record, and what happens when a change is
    accepted?
  \item Define mail merge and name its two parts.
\end{enumerate}

\subsection*{Solutions}
\begingroup\small
\begin{enumerate}[label=\arabic*.]
  \item Find: Ctrl+F. Find and Replace: Ctrl+H. Spelling and grammar: F7.
    Insert a hyperlink: Ctrl+K.
  \item A Page Break moves text to a new page within the same section. A
    Section Break starts a new section with its own layout. Only a Section
    Break allows different headers and footers.
  \item A bookmark is a named location or selection inside the document; it is
    a target and is not printed. A hyperlink is clickable text or an object
    that jumps to a target, which may be a bookmark, a file or a web address.
  \item Track Changes records every insertion, deletion and formatting edit,
    together with the author's name and the time (revision metadata). Accepting
    a change keeps it permanently and removes that change's revision metadata.
  \item Mail merge combines a main document (the standard letter or template)
    with a data source (a list of records such as names and addresses) to
    produce one personalised copy per record.
\end{enumerate}
\endgroup
\end{document}
